\documentclass[11pt,a4paper]{article}

% Idioma y fuentes (compilar con XeLaTeX)
\usepackage{fontspec}
\usepackage[spanish]{babel}

% Matemáticas, figuras y enlaces
\usepackage{amsmath, amssymb}
\usepackage{graphicx}
\usepackage[hidelinks]{hyperref}

\title{Título del artículo}
\author{Nombre Apellido}
\date{\today}

\begin{document}

\maketitle

\begin{abstract}
  Resumen breve del trabajo: el problema, el enfoque y los resultados principales.
\end{abstract}

\section{Introducción}
\label{sec:introduccion}

Este es un artículo de ejemplo. Las tildes, la ñ y los signos de apertura (¿?, ¡!) funcionan
directamente gracias a XeLaTeX. Para más detalles, véase la sección~\ref{sec:desarrollo} y la
referencia~\cite{knuth1984}.

\section{Desarrollo}
\label{sec:desarrollo}

Una ecuación numerada:
\begin{equation}
  \label{eq:euler}
  e^{i\pi} + 1 = 0
\end{equation}

y un sistema alineado:
\begin{align}
  \nabla \cdot \vec{E} &= \frac{\rho}{\varepsilon_0}, \\
  \nabla \cdot \vec{B} &= 0.
\end{align}

La ecuación~\eqref{eq:euler} relaciona cinco constantes fundamentales.

\section{Conclusiones}

Resume aquí las conclusiones principales.

\bibliographystyle{plain}
\bibliography{referencias}

\end{document}
